\section{Introduction}
\label{sec:introduction}

The question of whose population counts for redistricting has divided legal scholars, advocates, and political actors for generations. The most contested version of this question concerns noncitizens: approximately 22 million noncitizen adults reside in the United States, concentrated in urban metropolitan areas in a small number of high-immigration states. Whether those noncitizens should be counted for redistricting---and if so, whether they should count toward the total population or be excluded in favor of a citizen-only denominator---determines how political representation is distributed between urban immigrant communities and suburban citizen-majority communities.

The Supreme Court addressed one dimension of this question in \emph{Evenwel v. Abbott}~\citep{evenwel2016}: states \emph{may} use total population as their redistricting base, and the challengers' argument that the Equal Protection Clause \emph{requires} citizen VAP was rejected unanimously. But the Court conspicuously declined to hold that total population is the \emph{only} permissible base, leaving open whether states \emph{may} use citizen VAP instead. A forthcoming legal battle---perhaps already underway in some state capitals---will contest the other side of \emph{Evenwel}: whether the Constitution or federal statute \emph{requires} citizen VAP, or at minimum \emph{permits} it as an alternative.

This paper addresses four questions:
\begin{enumerate}
    \item What is the legal framework governing the choice between total population and citizen VAP, separately for congressional and state legislative redistricting?
    \item How large is the empirical divergence between total population and citizen VAP in the high-immigration states, and which districts are most affected?
    \item What is the partisan direction of switching from total population to citizen VAP, and how many congressional seats would plausibly shift?
    \item How does BISECT's \texttt{--population-source cvap} implementation enable systematic empirical analysis of the redistricting consequences?
\end{enumerate}

\subsection{Definitions}

\begin{itemize}
    \item \textbf{Total population}: The Census Bureau's enumerated total, including all persons regardless of citizenship or immigration status.
    \item \textbf{Total VAP}: Total voting-age population (persons 18 and older), including noncitizens.
    \item \textbf{Citizen VAP (CVAP)}: Citizen voting-age population from the Census Bureau's CVAP special tabulation, derived from American Community Survey estimates and published at the census-tract level.
    \item \textbf{Noncitizen VAP}: Total VAP minus citizen VAP; the population of voting-age noncitizens.
    \item \textbf{Noncitizen fraction}: Noncitizen VAP divided by total VAP in a given jurisdiction.
\end{itemize}

\subsection{Contributions}

\begin{itemize}
    \item Legal analysis of the congressional--state legislative distinction in population-base choice, including the Article I textual argument and the \emph{Evenwel} framework
    \item Empirical quantification of total-population--citizen-VAP divergence in California, Texas, New York, and Florida, at the congressional district and census-tract levels
    \item Partisan direction analysis showing consistently Republican-advantaging effect of citizen-VAP redistricting
    \item Estimate of 15--20 congressional seats affected by a nationwide shift to citizen-VAP redistricting
    \item Description of BISECT's CVAP implementation (\texttt{--population-source cvap}) and its data source (Census CVAP special tabulation)
    \item Analysis of \emph{Burns v. Richardson} and \emph{Evenwel} as the controlling legal anchors, with attention to their distinct constitutional logics
\end{itemize}
